% TOP_PROBLEM: {"id": 2736, "title": "Kirby Problem 1.77", "category_id": 7, "category": {"id": 7, "name": "topology", "display_name": "Topology", "description": "Properties preserved under continuous deformations.", "slug": "topology", "order_index": 7, "created_at": "2026-07-31T15:26:25.671Z"}, "status": "open", "problem_number": "KP-1.77", "set_id": 11}
% FIRST_POSTED: 2026-10-01
% ENTRY_KIND: statement_only
% UnsolvedMath ID 2736; source catalogue code KP-1.77
% !TeX program = lualatex
% Definition and sources only; this file is not an LLM solution attempt.
\documentclass[11pt,a4paper]{article}
\usepackage[margin=25mm]{geometry}
\usepackage{fontspec}
\setmainfont{lmroman10-regular.otf}[BoldFont=lmroman10-bold.otf,
 ItalicFont=lmroman10-italic.otf,BoldItalicFont=lmroman10-bolditalic.otf]
\usepackage{amsmath,amssymb,mathtools}
\usepackage{unicode-math}
\setmathfont{latinmodern-math.otf}
\AtBeginDocument{\let\setminus\smallsetminus}
\usepackage{xurl,hyperref}
\hypersetup{colorlinks=true,urlcolor=blue}
\setcounter{secnumdepth}{0}
\setlength{\parindent}{0pt}
\setlength{\parskip}{5pt}
\setlength{\emergencystretch}{3em}
\begin{document}
\section*{Kirby Problem 1.77}\label{2736}
{\small UnsolvedMath ID 2736 · KP-1.77 · Topology · Kirby's Problems in Low-Dimensional Topology}
\subsection{Definitions and mathematical statement}
Consider diagrams of knots in the plane, modulo planar isotopy. The three
Reidemeister moves change a curl, a canceling two-crossing pair, or the order of
three strands around a crossing. Let $d_R(D,E)$ be the smallest number of such
moves connecting diagrams $D,E$ of the same knot, assigning zero cost to planar
isotopy. Intermediate diagrams may have arbitrarily many crossings.
Write $c(D)$ for the number of crossings of $D$.

Determine sharp quantitative bounds, in terms of endpoint crossing numbers,
for $d_R(D,E)$. One precise worst-case formulation is
\[
R(n)=\max\{d_R(D,E):c(D)+c(E)\leq n,\ D,E\text{ represent the same knot}\}.
\]
The maximum uses combinatorial diagram types, of which there are finitely many
with bounded crossing number. Reidemeister's theorem makes each distance finite.
The problem asks for the growth of $R(n)$, or comparably sharp two-variable
bounds in $c(D),c(E)$, valid uniformly over knot types. A particularly strong
possibility is a polynomial upper bound. The quantity measures the number of
moves, rather than the greatest crossing number encountered during them.
\subsection{Short English statement}
How many Reidemeister moves are needed in the worst case to connect two diagrams of the same knot?
\subsection{Sources}
\begin{itemize}
\item[S1] ULAM.AI and the UnsolvedMath Contributors, supplied problems.json, record 2736 (KP-1.77), Kirby's Problems in Low-Dimensional Topology. Catalogue identity and source transcription; CC BY 4.0.
\url{https://huggingface.co/datasets/ulamai/UnsolvedMath}.
\item[S2] K3: A New Problem List in Low-Dimensional Topology, Problem 1.77. Expanded formulation of the supplied catalogue statement.
\url{https://math.berkeley.edu/sites/default/files/surv-295-ruberman-watermarked-author-pdf.pdf}.
\end{itemize}
\end{document}
